\section{AI Coding: después de que el código se abarata, la responsabilidad se encarece}

AI Coding es la línea más realista de este episodio. El anfitrión menciona que herramientas como Cloud Code hacen que la producción de código supere por mucho la velocidad humana, lo cual repercutirá en los reviewers de GitHub, los maintainers de código abierto y el code review interno de las empresas. Zhao Chenyang (赵晨阳) ya enfrenta dentro de SGLang la presión del código generado por IA, así que su respuesta es muy concreta.

Él no se opone a que la IA escriba código; al contrario, lo recibe con gusto. Lo esencial es esto: quien envía un PR debe poder explicar el design feature y estar dispuesto a responder por el resultado. Si alguien entrega un PR y ni siquiera puede aclarar los detalles ni la motivación del diseño, ese PR puede cerrarse sin ninguna contemplación. La IA no es un agente que asuma la responsabilidad por ti; si la usas para escribir código, cuando algo falle sigues siendo tú quien responde.

\begin{importantbox}{La regla central de los proyectos de código abierto frente a los PR de IA}
Se recibe con gusto que la IA escriba code, pero no que se haga sin pensar. El contribuyente debe explicar el diseño, entender los detalles y estar dispuesto a responder. Lo que la IA reduce es el costo de producir código, no la responsabilidad de mantener la calidad del proyecto.
\end{importantbox}

Este juicio tiene un sentido muy práctico. También se menciona el ejemplo de Amazon: se puede usar IA para escribir código, pero si algo sale mal hay que responder por ello. El mecanismo de responsabilidad termina por frenar el impulso de «generar 50K de código al día», porque cuanta más IA genera, más aumentan tanto la probabilidad de errores como el costo del review. Si un contribuyente solo busca acumular PR y registros de contribución, sin querer asumir la responsabilidad de mantenimiento a largo plazo del código, el proyecto de código abierto se desgastará hasta quebrarse.

El anfitrión plantea una observación muy importante: Git administra el resultado final del código, no el proceso de pensamiento. Vemos el diff final, pero no vemos cómo se formó ese diseño, cómo razonó el modelo, cómo interactuó el usuario con el coding agent, ni qué opciones se descartaron. Esto se conecta con el problema del debug log mencionado antes: lo verdaderamente escaso en la era de AI Coding no es solo el código en sí, sino el proceso de diseño, el nivel de pensamiento y la cadena de responsabilidad.

\begin{warningbox}{Git registra el código, no el juicio}
El gobierno del código en la era de la IA no puede basarse solo en el tamaño del diff, el número de PR o la cantidad de código. Dos PR con la misma cantidad de código pueden tener detrás un pensamiento, un riesgo y una responsabilidad de mantenimiento completamente distintos. Lo que en verdad habrá que registrar en el futuro es el decision flow, el proceso de review y las razones del diseño.
\end{warningbox}

Aquí se introduce la analogía entre el context graph y el SaaS. Lo verdaderamente valioso de un SaaS no es solo el materialized data, sino también cómo entraron esos datos al sistema, el proceso de negocio detrás y las offline discussion. Con Coding ocurre algo similar: el código final es solo el resultado; el flujo de diseño, el proceso de interacción y el contexto detrás son el activo más profundo. Quien logre capturar estos procesos podría formar, en la era de AI Coding, un nuevo ciclo cerrado de datos y capacidades.

\subsection{Resumen de la sección}

AI Coding abarata la producción de código, pero vuelve más importantes la responsabilidad, el diseño y el review. Los proyectos de código abierto no pueden resolver el problema prohibiendo el código de IA; deben establecer límites de responsabilidad: el contribuyente debe entender, explicar y mantener su propio PR. Los datos clave de la colaboración en código en el futuro tal vez no sean solo el diff final, sino el decision flow, el context graph y las razones del diseño.
